\section{Pregunta 5: El Kitty Problem y el Principio Abierto/Cerrado}

% Responsable: Miguel Andrés Cazorla Zamora

\subsection{Descripción del problema}

El \textit{Kitty Problem} parte de un sistema agregador de taxis dividido en
servicios: \texttt{TaxiUI} (atiende al cliente), \texttt{TaxiFinder} (busca taxis
disponibles), \texttt{TaxiSelector} (elige el más adecuado según coste, tiempo o
lujo), \texttt{TaxiDispatcher} (ordena el servicio) y \texttt{TaxiSupervisor}
(controla el viaje) \parencite{martin2017}. Los servicios se despliegan por
separado y parecen desacoplados.

Un día se pide un requisito nuevo: la entrega de gatitos a domicilio. Hay
conductores alérgicos a los gatos que deben excluirse, vehículos no aptos e
instrucciones de recogida y devolución propias. El resultado es que \emph{todos}
los servicios deben modificarse, y esos cambios deben coordinarse y desplegarse
juntos. Los servicios estaban desacoplados en ejecución, pero no frente al
cambio \parencite{martin2017}.

\subsection{¿Por qué la división por verbos provocó el problema?}

Los servicios \texttt{TaxiFinder} o \texttt{TaxiSelector} representan
\emph{pasos del algoritmo} de reserva, no conceptos de negocio. Es una
descomposición funcional: el concepto \enquote{tipo de viaje} queda repartido
entre todos los servicios, de modo que cuando ese concepto cambia (aparecen los
gatitos), el cambio se propaga a todos ellos. \textcite{richardsford2020} lo
describen como una partición técnica frente a una partición por dominio: en la
primera, un cambio de dominio atraviesa todas las particiones; en la segunda,
queda confinado a una sola. Además, todos los servicios comparten la misma
representación del viaje, por lo que siguen acoplados a través de los datos.

\subsection{Diseño basado en DDD y Onion/Clean}

DDD propone dividir el sistema según los conceptos del dominio y sus ejes de
cambio. Aquí la observación clave es que la entrega de gatitos \emph{es un tipo
de viaje}: comparte el flujo general (buscar, seleccionar, despachar,
supervisar) pero varía en reglas concretas. Por tanto se separan:

\begin{itemize}[nosep]
  \item \textbf{\texttt{Rides}}: núcleo con la lógica estable del viaje y los
        puntos de extensión (clases abstractas e interfaces como \texttt{Ride} o
        \texttt{DriverEligibility}).
  \item \textbf{\texttt{Kittens}}: plugin que implementa esas abstracciones con
        las reglas propias de los gatitos (p.\,ej.\ excluir conductores
        alérgicos).
\end{itemize}

Siguiendo la Regla de Dependencia \parencite{martin2017} y la Arquitectura Onion
\parencite{palermo2008}, \texttt{Rides} ocupa el centro y \texttt{Kittens}
depende de él, nunca al revés (Figura~\ref{fig:p5-plugins}). Cada servicio puede
mantener su despliegue, pero internamente opera sobre las abstracciones de
\texttt{Rides} mediante polimorfismo (\textit{Template Method} o
\textit{Strategy}), sin conocer el tipo concreto de viaje \parencite{martin2017}.
Así, añadir gatitos solo afecta a un componente: \texttt{Kittens}.

\begin{figure}[ht]
\centering
\begin{tikzpicture}[
    font=\small,
    box/.style={draw, rounded corners=2pt, minimum width=3.3cm, minimum height=0.8cm, align=center},
    impl/.style={-{Triangle[open,length=2.4mm]}, thick, dashed},
  ]
  \node[box, fill=blue!8] (ride) {\texttt{Ride}\\\footnotesize(abstracta)};
  \node[box, fill=blue!8, right=0.8cm of ride] (crit) {\texttt{DriverEligibility}\\\footnotesize(interfaz)};
  \node[box, fill=orange!12, below=1.6cm of ride] (kride) {\texttt{KittenRide}};
  \node[box, fill=orange!12, below=1.6cm of crit] (kelig) {\texttt{NoCatAllergy}\\\texttt{Eligibility}};
  \begin{scope}[on background layer]
    \node[draw, thick, rounded corners, fill=blue!3, inner sep=8pt, fit=(ride)(crit),
          label={[font=\bfseries\small]west:\texttt{Rides}}] {};
    \node[draw, thick, rounded corners, fill=orange!4, inner sep=8pt, fit=(kride)(kelig),
          label={[font=\bfseries\small]west:\texttt{Kittens}}] {};
  \end{scope}
  \draw[impl] (kride) -- (ride);
  \draw[impl] (kelig) -- (crit);
\end{tikzpicture}
\caption{El plugin \texttt{Kittens} implementa las abstracciones del núcleo
\texttt{Rides}. Las dependencias apuntan hacia el núcleo.}
\label{fig:p5-plugins}
\end{figure}

\subsection{Principio de Abierto/Cerrado y papel de los Adaptadores}

El OCP establece que un artefacto debe estar abierto a la extensión y cerrado a
la modificación \parencite{martin2017}. En el diseño original, añadir gatitos
exigía modificar cinco servicios. En el diseño propuesto, el núcleo
\texttt{Rides} y los servicios existentes no se tocan (cerrados a la
modificación) y el nuevo comportamiento se añade creando el componente
\texttt{Kittens} (abierto a la extensión).

Los Adaptadores hacen posible esta extensión. En términos de la Arquitectura
Hexagonal \parencite{cockburn2005}, las abstracciones de \texttt{Rides} son
\emph{puertos} y \texttt{Kittens} es un \emph{adaptador} que traduce las reglas
de los gatitos al contrato que entiende el núcleo, que nunca conoce el concepto
\enquote{gatito}. Si los gatitos provinieran de un sistema externo, un adaptador
de servicio encapsularía esa integración tras un puerto del dominio
\parencite{delatorre2010}. Finalmente, el componente \textit{Main} conecta el
plugin mediante factorías o inyección de dependencias \parencite{martin2017},
por lo que incorporar un nuevo tipo de viaje se reduce a registrar un nuevo
adaptador.
